\documentclass[../../main.tex]{subfiles}
\begin{document}

\begin{center}
  {\fontsize{48pt}{58pt}\selectfont \textbf{2000}\normalsize\textbf{年度}}
\end{center}

\begin{table}[htb]
  \centering
  \caption{出題テーマと難易度}
  \begin{tabular}{|c|p{9cm}|c|}
    \hline
    \textbf{問題番号} & \textbf{テーマ} & \textbf{難易度} \\
    \hline
    第1問 & 直角二等辺三角形に内接する楕円の面積 & \\
    \hline
    第2問 & 複素数平面上の直線と円 & \\
    \hline
    第3問 & 漸化式で定まる関数の極限とグラフ & \\
    \hline
    第4問 & 円周上の点と動く線分の衝突 & \\
    \hline
    第5問 & 十進表示の桁の条件と数え上げ & \\
    \hline
    第6問 & 行列の積と立体の体積 & \\
    \hline
  \end{tabular}
\end{table}

\begin{table}[htb]
  \centering
  \caption{解答}
  \renewcommand{\arraystretch}{1.4}
  \begin{tabular}{|c|c|p{9cm}|}
    \hline
    \textbf{問題番号} & \textbf{小問} & \textbf{答え} \\
    \hline
    第1問 & -- & 面積の最大値は$\dfrac{\sqrt{3}}{9}\pi$ \\
    \hline
    第2問 & -- & $w=\alpha\beta$と$P,Q$が指定された円周上にあることは
    同値（証明問題） \\
    \hline
    第3問 & (1) & $p_k=1+\dfrac{1-c}{c}(1-h)^k$ \\
    \cline{2-3}
    & (2) & $g(a)=\dfrac{c}{c+(1-c)e^{-a}}$ \\
    \cline{2-3}
    & (3) & $c=2,1,\dfrac{1}{4}$のグラフを図示（解答参照） \\
    \hline
    第4問 & (1) & $0<v<\dfrac{3}{2\pi}$，または
    $v>\dfrac{9}{4\pi}$ \\
    \cline{2-3}
    & (2) & $\dfrac{9}{28\pi}<v\le\dfrac{9}{16\pi}$，
    または$\dfrac{3}{2\pi}\le v\le\dfrac{9}{4\pi}$ \\
    \hline
    第5問 & (1) & $1728$個 \\
    \cline{2-3}
    & (2) & $8695$ \\
    \hline
    第6問 & (1)〜(3) & \\
    \hline
  \end{tabular}
\end{table}

\paragraph{試験形態}
6問．試験時間は要確認．

\end{document}
